\section{Learned Models}
\label{sec:models}

\subsection{Why transformers, and which}
The rule has two structural limits that no amount of tuning removes. It never looks further than one sentence
from a citation, so it finds 156 of the 231 near gold sentences, and it finds \textbf{none of the 112 far
ones}. Going further requires judging meaning: does this sentence in the Discussion restate the result that
the caption of Figure~2 describes? That is a job for a pre-trained transformer \emph{encoder}, a model that
reads text and produces a score, as opposed to a generative model that writes text.

There is a second reason to move beyond the rule. \bof{} and its filters are heuristics derived from
biomedical articles: the position of implicit sentences, the wording of supplementary-material mentions, the
names of Methods sections. Other fields write about their figures differently, and such rules need not
transfer. A model that judges whether a sentence and a caption are about the same thing, from their meaning,
is a more trustworthy basis for papers from other domains.

We use two encoders, in two different roles.
\begin{itemize}
  \item \textbf{DeBERTa-v3} \citep{he2023debertav3} as a \emph{cross-encoder}: one model reads the caption and
        a candidate sentence together and outputs one score. DeBERTa keeps separate representations for what a
        token is and where it is, and is among the strongest open encoders for judging a pair of texts. Its
        512-token limit is sufficient for a caption and a sentence. We use the base model (184M parameters)
        and then the large one (435M).
  \item \textbf{ModernBERT-large} \citep{warner2024modernbert} (395M parameters) as a \emph{span tagger}: one
        model reads a long stretch of the paper and outputs a score for every sentence in it. Its
        8{,}192-token context makes this possible on a single 15\,GB GPU.
\end{itemize}
The two see the problem differently. The cross-encoder judges one sentence against the caption and is told the
sentence's position as text; the span tagger sees the neighbouring sentences and the flow of the paragraph
directly. As Section~\ref{sec:ensemble} shows, their errors differ enough that averaging them helps. We did not
use generative language models: small ones were weak in our tests (Section~\ref{sec:othertrials}) and large
ones conflict with the requirement to run locally. All training ran on NVIDIA T4 GPUs on Kaggle.

\subsection{The data problem}
All learning has to come from \textbf{549 gold sentences in 15 papers}, of which only \textbf{112 are far}.
The training pool contains 11{,}998 (figure, sentence) pairs, so 4.6\% are positive. We searched for further
labelled data and found none: the large figure datasets label explicit mentions only
(Section~\ref{sec:related}). A model with 400 million parameters will memorise 15 papers. We therefore tried
three ways of creating more training signal: free labels from other papers, a two-stage curriculum, and
augmentation of the real data.

\subsection{Manufactured labels: two ideas that did not transfer}
\label{sec:fakelabels}

\paragraph{Free labels from citations (``M2'').} An implicit reference is, loosely, an explicit reference
without the citation. So we manufactured them: in 6{,}049 open-access papers we took every sentence that cites
a figure in parentheses, deleted the citation, and used the result as a positive example for that figure's
caption. Negatives were sentences that cite a different figure, or that occur before the figure is first
cited. This gave about 113{,}000 pairs, on which we trained a DeBERTa-v3-base cross-encoder. On held-out
papers of the same kind the model was good (AUC 0.873). On \figsum{} it was not: for near sentences its AUC
was 0.614, \emph{below} plain TF-IDF similarity to the caption (0.639). Fine-tuning it further on \figsum{}
pairs only rediscovered the \pmk{1} rule, and starting from the M2 weights was worse than starting from the
plain pre-trained model.

\paragraph{A two-stage curriculum.} The second idea was to separate two skills. Stage~1 learns to match
captions and sentences from \emph{all} gold sentences, with every figure number masked (``Fig.~[X]'') so that
the model must use content. Stage~2 starts from stage~1 and learns implicit references given the explicit
ones. Stage~1 did not learn anything that transfers. In the first attempt (all negatives, a positive weight of
about 20) the loss stayed flat and the near AUC on held-out papers was 0.521. In the second attempt (four
negatives per positive, resampled each epoch, a lower learning rate) the training loss did fall, but the
held-out near AUC was still 0.52, ranging from 0.38 to 0.61 across folds: the model memorised its training
papers. Stage-2 models scored F1 0.666 to 0.682, well below the simple cross-encoder (0.728). We closed the
idea after two attempts.

\paragraph{Why both failed.} A sentence written to carry ``(Fig.~$N$)'' is a statement of a result. The
implicit sentences that humans mark are often of a different kind: the sentence that sets up the experiment
(``To test whether \ldots''), or a sentence in the Discussion that interprets the result. Manufactured
positives teach the model to recognise result statements, which is not the distinction the task needs. The
``figure-talk'' classifier of Appendix~\ref{app:other} failed for the same reason. One part of the
two-stage intuition did survive: training on the explicit sentences as well, with figure numbers
\emph{visible}, clearly helps, as the next subsection shows.

\subsection{What worked: every real pair, lightly augmented}
\label{sec:dataset}
The simplest recipe was the best. We train on every (figure, sentence) pair of the training papers, including
the explicit sentences, and add three light techniques.
\begin{itemize}
  \item \textbf{Citation-deleted copies (``A3'').} For a sentence that cites $F$, we also train on a copy from
        which the citation has been deleted and which is presented as not citing $F$; its label is unchanged.
        This stops the model from relying on the citation marker alone. In the stage-2 model, removing A3 cost
        0.038 F1 (0.670 to 0.632).
  \item \textbf{Paraphrases by back-translation} \citep{sennrich2016backtranslation}. Sentences are translated
        to German and back with MarianMT \citep{junczys2018marian}, with figure mentions protected, giving
        paraphrases of 1{,}773 of the 2{,}232 training sentences. During training a sentence is swapped for
        its paraphrase with probability 0.3. Together with A3 this raised the far AUC of the base cross-encoder
        from 0.674 to 0.722 at unchanged F1.
  \item \textbf{Sampling of far negatives.} For a given figure almost every far sentence is a negative. Each
        epoch uses all gold pairs, all explicit and near pairs, and a freshly drawn 20\% of the far non-gold
        pairs. This brings the ratio of negatives to positives from about 20:1 to about 5:1 and makes an epoch
        three to four times shorter. We did not run a control without sampling, so we claim no accuracy gain
        for it.
\end{itemize}
The final dataset file holds 18 papers, 2{,}695 sentences with their paragraph, section and citations,
14{,}792 (figure, sentence) pairs (11{,}998 for training), 657 gold labels, the paraphrases, 125
citation-deleted copies, a masked copy of every sentence, and the fold of every paper.

\subsection{C0: the pairwise cross-encoder}
\label{sec:c0}
\paragraph{Input.} For a figure $F$ and a candidate sentence, the model receives two texts. Text~A is the
caption, with a leading ``Fig.~$N$.'' removed and truncated to 700 characters. Text~B is the candidate,
preceded by a tag written in plain words and by the nearest sentence that cites $F$:

\begin{Verbatim}[fontsize=\footnotesize,frame=single,framesep=6pt,rulecolor=\color{navy!30}]
A: Production of interleukin (IL)-17 by peripheral blood mononuclear cells ...
B: [cites this figure: no | 1 sentence before a citation of this figure |
    same paragraph | section: Results | cites: none]
   Citing: We detected an increase in IL-17 level after stimulation ... (Fig. 2a).
   Candidate: To evaluate the effects of inflammatory cytokines released by ...
\end{Verbatim}

The tag has five fields: whether the candidate itself cites $F$; its distance and direction to the nearest
citation of $F$ (or ``this is a citation'', or ``no citation of this figure in the paper''); whether it shares
a paragraph with that citation; its top-level section; and what else it cites (other figures, a table,
supplementary material, or nothing). The tag gives the model everything the rule knows, as text that it can
learn to weigh, in place of a hard distance limit. The citing sentence is included because implicit sentences
often continue what the explicit one says.

\paragraph{Model.} The pair is tokenised to at most 320 tokens and passed through DeBERTa-v3. The final hidden
state of the first token goes through the model's pooling layer and a linear layer to a single logit; its
sigmoid is the probability that the sentence is associated with $F$. Every sentence of the paper is scored
for every figure.

\paragraph{Training.} Binary cross-entropy, with positives weighted by the ratio of negatives to positives
(about 5); AdamW \citep{loshchilov2019adamw} with learning rate $2\times10^{-5}$ and weight decay 0.01; 10\%
linear warm-up followed by linear decay; batch size 16; three epochs; mixed precision. In a representative
fold the training loss per epoch was 0.97, 0.63, 0.57. Full hyperparameters are in
Appendix~\ref{app:hyper}.

\paragraph{Result.} C0 with the base model scores F1 0.728 out-of-fold, below \bof{} (0.752), although its
ranking is good (near AUC 0.822). An earlier version of the same model on the 13-paper protocol had shown
the same pattern (F1 0.721 against 0.767 for \bof{}, near AUC 0.847). The diagnosis that guided the following
steps was therefore: \emph{the model ranks sentences well; the gain is lost when scores are turned into a set}.

\subsection{Bigger model and the first fix to the decision step}
\textbf{C0L} is the same model with DeBERTa-v3-large (learning rate $10^{-5}$; micro-batches of 4 accumulated
to 16 and gradient checkpointing, since a batch of 8 did not fit in memory). Its near AUC rose to 0.848, the
best so far, yet its F1 was 0.725.

The cause was the decision step. We had been selecting a sentence when the \emph{z-score of its logit within
its paper} exceeded a threshold, on the reasoning that papers differ in scale. But standardising within a
paper discards how confident the model is in absolute terms, and the large model is confident. Thresholding
the raw logit instead gave F1 \textbf{0.752}, level with the rule.
From this point on, the decision step received as much attention as the model.

Two further variants of the large cross-encoder did not help: adding the augmentations (C0aL, 0.723) and
appending the sentences immediately before and after the candidate to the input (C0Lctx, 0.735, with a lower
far AUC). The second result suggested that context should not be bolted on to a pairwise model but read
natively.

\subsection{SPAN: the whole-paper span tagger}
\label{sec:span}
\paragraph{Input.} For a figure $F$, the input begins with the caption and continues with the sentences of the
paper in order. Each sentence is preceded by a marker token (we reuse the tokenizer's \texttt{[MASK]} token
purely as a placeholder), and by short inline annotations of structure:

\begin{Verbatim}[fontsize=\footnotesize,frame=single,framesep=6pt,rulecolor=\color{navy!30},commandchars=\\\{\}]
[CLS] Figure 2. Production of interleukin (IL)-17 by ... [SEP]
[MASK] \S{} Results \P{} To evaluate the effects of inflammatory cytokines ...
[MASK] [cites this figure] We detected an increase in IL-17 level ...
[MASK] [cites Figure 3] ...
   ...
[MASK] \S{} Discussion \P{} ...
\end{Verbatim}

``\S{} \emph{section}'' appears when the top-level section changes, ``\P'' at the start of each paragraph,
``[cites this figure]'' on sentences that name $F$, and ``[cites Figure $g$]'' on sentences that name other
figures. The caption header is limited to 200 tokens, each sentence to 96 tokens, and a window to 2{,}048
tokens.

\paragraph{Windows.} A paper is longer than one window, so windows tile the paper: a window is filled greedily
with whole sentences, and the next window starts two-thirds of the way through the previous one, so that
consecutive windows overlap by a third. A sentence near the edge of one window is thus seen again with context
on both sides. A sentence that falls in two windows receives the mean of its two logits.

\paragraph{Model.} ModernBERT-large encodes the window once. A linear layer ($1024\rightarrow1$) applied to the
hidden state at each marker gives one logit per sentence. A single forward pass therefore scores every sentence
in the window, where the cross-encoder needs one pass per sentence, and each score is computed with the
neighbouring sentences, paragraph boundaries and other figures' citations in view. This lets the model learn
what the CRF of Appendix~\ref{app:other} was meant to capture, namely that annotators mark runs.

\paragraph{Training.} For each figure we keep every window that contains a citation or a gold sentence of that
figure, and 30\% of the remaining windows. Citation-deleted copies and paraphrases are each applied with
probability 0.3. The loss is binary cross-entropy over all markers, with positives weighted by the
negative-to-positive ratio capped at 8; learning rate $2\times10^{-5}$; effective batch of 8 windows; five
epochs; mixed precision and gradient checkpointing.

\paragraph{Result.} SPAN has the best ranking of any single model (near AUC 0.872, far AUC 0.701) and an F1 of
0.747 with a raw threshold.

\subsection{Anchoring: a decision rule calibrated per figure}
\label{sec:anchor}
A single threshold has to serve every figure, but the models score some figures systematically higher than
others, for instance figures with long and specific captions. A figure's own citing sentences offer a natural
reference point: they are known to be about the figure. We define the \textbf{anchored score} of a sentence as
\[
a(F,s)=\text{logit}(F,s)-\operatorname{mean}_{e\in E(F)}\text{logit}(F,e),
\]
where $E(F)$ is the set of sentences that cite $F$ (if a figure is never cited, we subtract the median logit
of its sentences). The prediction is $E(F)\cup\{s: a(F,s)\ge t\}$ with one threshold $t$ for all figures. In
words: select a sentence if the model scores it nearly as high as it scores the figure's own citations.

For example, suppose that for some figure the citing sentences score about $-2$ and a true implicit sentence
scores $-4$. A raw threshold of 0 rejects it; the anchored score is $-2$ and it is selected. Table~\ref{tab:decision}
compares the three decision rules on the same C0L scores.

\begin{table}[h]
\centering
\small
\begin{tabular}{lc}
\toprule
Decision rule (C0L, explicit sentences always included) & OOF F1\\
\midrule
z-score of the logit within the paper $\ge t$ & 0.725\\
Raw logit $\ge t$ & 0.752\\
Anchored score $\ge t$ & \best{0.756}\\
\bottomrule
\end{tabular}
\caption{The same model scores under three decision rules. Thresholds are chosen out-of-fold.}
\label{tab:decision}
\end{table}

\subsection{Ensemble}
\label{sec:ensemble}
The ensemble score of a pair is the mean of the C0L logit and the SPAN logit; it is then anchored as above.
The ensemble ranks better than either model (near AUC 0.892 against 0.848 and 0.872; far AUC 0.783) and
scores \textbf{F1 0.766} out-of-fold, $+0.014$ over \bof{}. This
configuration, with the threshold $t=-3.6$ that maximises F1 on the out-of-fold scores of the training papers,
was registered as the primary model for the sealed test.

\subsection{Early stopping}
After the test (Section~\ref{sec:test}) we returned to the span tagger with validation-based early stopping.
Training runs for up to eight epochs. Every half epoch the model is scored on validation papers; the metric is
average precision, smoothed over the last two checks; training never stops before epoch 3 and stops after two
epochs without a new best, restoring the best weights. The median best epoch was 4.25, and 5 of the 16 jobs ran
the full eight epochs. The resulting model, \textbf{SPAN\_ES}, improves the single span tagger from 0.744 to
0.763 (anchored). Because the validation papers are the three former test papers, models from this point on
have no clean test number.

\subsection{Text from the figure image}
\label{sec:vis}
All the models above know a figure only through its caption. We tested whether text taken from the image
itself helps. Figures were cropped from the PDFs by locating each caption and taking the image or vector
region above it; this succeeded for 57 of the 91 figures. For each crop we produced (i) the text visible in
the figure, from the PDF text layer for vector figures and otherwise by prompting a small vision--language
model, Qwen2-VL-2B \citep{wang2024qwen2vl}, to list the labels it can read, and (ii) a literal description
from the same model (chart type, axes, legend, extremes). Both were appended to the caption header of the
span tagger (\textbf{SPAN\_VIS}); figures without a crop received a note saying that no image was available.

The result was negative: SPAN\_VIS scores 0.742 against 0.763 for the otherwise identical SPAN\_ES. We see
three likely reasons. Biomedical captions are long and already describe what the figure shows. The small
general-purpose vision model was often wrong; it described a table as a bar chart, for example. And a third of
the figures had no crop at all.

This agrees in part with the figure-captioning literature. \citet{yang2023scicapplus} report that generic
image features add nothing once the mentioning text is available, and \citet{huang2023summaries} that most of
the information needed for a caption is in the surrounding text. It does \emph{not} agree with them in one
respect: in SciCap+, OCR text did improve captioning. That work concerns computer-science figures, which carry
more text (axis labels, legends) and have shorter captions, and a generation task; so our negative result
should be read as specific to biomedical papers and a small vision model, and OCR remains worth testing on
other kinds of papers. A three-model ensemble that includes SPAN\_VIS does reach the best out-of-fold F1 we
observed (0.776), which suggests that the visual variant makes somewhat different errors even though it is not
better on its own.

\subsection{Summary of the progression}
Table~\ref{tab:progress} and Figure~\ref{fig:progress} list the steps in the order we took them.

\begin{table}[h]
\centering
\small
\begin{tabularx}{\linewidth}{rLrrr}
\toprule
\# & Step & OOF F1 & Near AUC & Far AUC\\
\midrule
0 & Rule \bof{} & 0.752 & -- & --\\
1 & C0: DeBERTa-v3-base cross-encoder, every pair (z-score rule) & 0.728 & 0.822 & 0.674\\
2 & Two-stage curriculum (stage 2 from stage 1) & 0.670 & 0.772 & 0.763\\
3 & C0a: C0 + citation-deleted copies + paraphrases & 0.728 & 0.833 & 0.722\\
4 & C0L: DeBERTa-v3-large & 0.725 & 0.848 & 0.653\\
5 & C0L with a raw-logit decision rule & 0.752 & 0.848 & 0.653\\
6 & C0Lctx: C0L + neighbouring sentences in the input & 0.735 & 0.832 & 0.621\\
7 & SPAN: ModernBERT-large span tagger & 0.747 & 0.872 & 0.701\\
8 & Anchored decision rule (C0L alone) & 0.756 & 0.848 & 0.653\\
\rowcolor{hl}
9 & Ensemble C0L + SPAN, anchored \emph{(pre-registered)} & \best{0.766} & 0.892 & 0.783\\
10 & SPAN\_ES: span tagger with early stopping & 0.763 & 0.882 & 0.785\\
11 & SPAN\_VIS: SPAN\_ES + figure OCR and description & 0.742 & 0.882 & 0.764\\
12 & Ensemble C0L + SPAN\_ES + SPAN\_VIS & \best{0.776} & 0.899 & 0.778\\
\bottomrule
\end{tabularx}
\caption{Out-of-fold results on the 15 training papers, in chronological order. Explicit sentences are always
included in the selection. Steps 10--12 used the former test papers for early stopping.}
\label{tab:progress}
\end{table}

\begin{figure}[h]
\centering
\begin{tikzpicture}
\begin{axis}[
  width=0.95\linewidth, height=5.6cm, ymin=0.65, ymax=0.80,
  xmin=0.5, xmax=12.5, xtick={1,...,12}, xlabel={Step (numbering of Table~\ref{tab:progress})},
  ylabel={Out-of-fold F1}, ymajorgrids, grid style={soft}, axis lines*=left,
  tick label style={font=\small}, label style={font=\small},
  legend style={font=\small, at={(0.98,0.04)}, anchor=south east, draw=none, fill=none}]
\addplot[accent, thick, dashed, domain=0.5:12.5, samples=2] {0.752};
\addlegendentry{Rule \bof{} (0.752)}
\addplot[steel, thick, mark=*, mark size=2pt] coordinates
 {(1,0.728) (2,0.670) (3,0.728) (4,0.725) (5,0.752) (6,0.735) (7,0.747) (8,0.756) (9,0.766) (10,0.763) (11,0.742) (12,0.776)};
\addlegendentry{Learned models}
\end{axis}
\end{tikzpicture}
\caption{Out-of-fold F1 at each step. The learned models pass the rule only after the decision step is fixed
(steps 5 and 8) and the two models are combined (step 9).}
\label{fig:progress}
\end{figure}

\subsection{Controlling overfitting}
With 15 training papers, overfitting is the default outcome; stage~1 of the curriculum showed it plainly
(training loss falling, held-out AUC at chance). Our controls were: reporting only out-of-fold numbers, so
that memorisation cannot inflate a result; few epochs with a small learning rate, weight decay and a decaying
schedule; augmentation and fresh sampling of negatives in every epoch; averaging three seeds; and, in the last
run, early stopping on validation papers. The three epochs of the cross-encoder follow common practice for
fine-tuning encoders of this family \citep{devlin2019bert} and were not tuned; the early-stopping run later
confirmed that the five epochs chosen for the span tagger were close to the median optimum of 4.25.
